\section{The exact constructive theory and the descent principle}
\label{sec:foundations}

Throughout, $\N=\{0,1,2,\ldots\}$ and $\Z$ is a coded set of signed natural
numbers. For $n\geq 3$, write
\[
 \FLT(n):\quad
 \forall x,y,z\in\N\;\bigl(x>0\wedge y>0\wedge z>0
       \longrightarrow x^n+y^n\ne z^n\bigr).
\]
The main assertions concern unrestricted sizes of $x,y,z$. No finite search
over possible counterexamples is being substituted for this assertion.

\subsection{Language, axioms, and logical rules}

Our $\ipra$ is first-order intuitionistic primitive recursive arithmetic.
Its numerical language has $0$, successor $S$, and a function symbol for
each specified primitive recursive construction, of every finite arity.
Initial functions are zero, successor, and projections; functions are closed
under composition and primitive recursion. For a recursion symbol $R$,
with recursion argument first, the defining equations are
\[
 R(0,\vec a)=f(\vec a),\qquad
 R(Sn,\vec a)=g(n,R(n,\vec a),\vec a).
\]
The other defining equations are those of projections and composition.
The arithmetic axioms also include $Sx\ne0$, injectivity of $S$, and
induction for every quantifier-free formula $A(n,\vec a)$:
\[
 A(0,\vec a)\ \longrightarrow\
 \bigl(\forall n\,(A(n,\vec a)\longrightarrow A(Sn,\vec a))\bigr)
 \ \longrightarrow\ \forall n\,A(n,\vec a).
\]
The logical language has equality, $\bot$, $\wedge$, $\vee$, $\to$,
$\forall$, and $\exists$. Its rules are the usual intuitionistic introduction
and elimination rules, assumption, equality reflexivity and substitution,
and substitution instances of arithmetic axioms. Quantifier introduction
and existential elimination have their usual fresh-variable conditions.
Negation abbreviates implication to $\bot$. There is no unrestricted law of
excluded middle, double-negation elimination, choice principle, or
induction for arbitrary quantified formulas.

This is precisely the presentation documented by the types
\texttt{PR}, \texttt{PTm}, \texttt{PRAAx}, and \texttt{Prf} in the inspected
ProveIt sources \cite{proveit}. In particular, \texttt{PRA.qfInd} is the
appropriate formal rule for the inductions below. The repository's
\texttt{pra\_iff\_native} translates an already constructed derivation to
the native scalar calculus. It does not turn an arbitrary theorem about
Lean's natural numbers into an $\ipra$ derivation.

\subsection{Decidable notation and bounded constructions}

For readability we use subtraction of integers, division with remainder,
gcd, divisibility, order, finite lists, and bounded quantifiers. These are
definitional conveniences. Integers can be encoded by $2a$ for $a\geq0$
and $2a-1$ for $-a<0$. A pair can be encoded by
$(a+b)^2+a$; both components of a valid code are bounded by the code, so
decoding by bounded search is primitive recursive. Repeated pairing codes
fixed tuples. Length and entries of a finite list can instead be coded by
a length and a base-$B$ integer with $B$ larger than every entry.

Every bounded loop here has a numerical clock fixed before the loop begins.
The loop state is a tuple, its update is primitive recursive, and its final
state is defined by primitive recursion on that clock. If an algorithm
stops early, subsequent updates leave its state unchanged. Searches return
the first satisfactory candidate and a success bit, or a fixed default.
An existence proof is always accompanied by a bound that makes the success
bit equal to one.

A primitive recursive predicate $P$ is represented by a characteristic
function $\chi_P\in\{0,1\}$. For example, a bounded universal quantifier is
represented by the finite product of its characteristic values, and a
bounded existential quantifier by a bounded scan. Induction on a formula
containing such notation means induction on the corresponding
quantifier-free formula involving these characteristic functions. The
equivalences with the bounded notation follow from the same finite-product
or scan recursions. Thus decidable case splits used below are legitimate
intuitionistically.

\begin{remark}[A necessary distinction]
We are not claiming that every true primitive recursive identity is
provable in $\ipra$, or that an arbitrary terminating algorithm is primitive
recursive. We give concrete bounded constructions and their arithmetic
invariants. Nor do we assume a primitive recursive universal evaluator for
all primitive recursive function codes; there is no such evaluator. Each
fixed construction used in a proof is expanded into its own function
symbols and defining equations.
\end{remark}

\subsection{Finite descent, proved using only quantifier-free induction}

\begin{lemma}[Clocked descent]\label{lem:clocked}
Let $P(s)$ be a primitive recursive predicate on coded states, and let
$D(s)$ and $r(s)$ be primitive recursive functions. Suppose that, by
quantifier-free arithmetic and decidable case distinctions, one has proved
\[
 P(s)\longrightarrow
 \bigl(P(D(s))\wedge r(D(s))<r(s)\bigr).
\tag{D}
\]
Then $\ipra$ proves $\forall s\,\neg P(s)$.
\end{lemma}

\begin{proof}
Define $s_0=s$ and $s_{k+1}=D(s_k)$ by primitive recursion. Fix a putative
state $s$ satisfying $P(s)$, and put $R=r(s)$. Induct on $k$ with the
quantifier-free invariant
\[
 k\leq R+1\ \longrightarrow\
       \bigl(P(s_k)\wedge r(s_k)+k\leq R\bigr).
\tag{I}
\]
The case $k=0$ is immediate. If $k+1\leq R+1$, the inductive hypothesis
applies at $k$, and (D) gives
$r(s_{k+1})+1\leq r(s_k)$. Adding $k$ proves the next case of (I),
including preservation of $P$. At $k=R+1$ we obtain
$r(s_{R+1})+R+1\leq R$, impossible for natural numbers. This is a
derivation of $P(s)\to\bot$, followed by universal introduction. No
infinite sequence and no least counterexample has been postulated.
\end{proof}

One can equally use a finite search for a least counterexample inside a
box whose size is supplied by an alleged solution. The clocked version
makes the induction formula and the computation bound more visible, so it
will be used throughout.

\subsection{Elementary integer arithmetic needed by the descents}
\label{sec:integer-tools}

Here are the elementary tools, including their constructive justifications.
They will also fix what is meant by ``extract a coprime power'' later.

\begin{lemma}[Division, gcd, and prime powers]\label{lem:integer-tools}
The following operations and facts are available in $\ipra$.
\begin{enumerate}
 \item For $b>0$, division with remainder produces $a=bq+r$ with
 $0\leq r<b$. Euclid's algorithm produces a gcd and integer Bezout
 coefficients.
 \item A positive integer can be factored into primes by bounded trial
 division. If a prime $\ell$ divides $uv$, then it divides $u$ or $v$.
 Valuations of nonzero integers obey
 $v_\ell(uv)=v_\ell(u)+v_\ell(v)$.
 \item If $m\geq1$ and pairwise coprime nonzero integers $a_1,\ldots,a_t$ have product
 an $m$th power, then each $|a_i|$ is an $m$th power. When $m$ is odd,
 each signed $a_i$ is an $m$th power. If $a,b>0$ are coprime and
 $ab=2u^2$, the odd factor is a square and the even factor is twice a
 square.
\end{enumerate}
\end{lemma}

\begin{proof}
The quotient is the largest $q\leq a$ with $bq\leq a$; its existence and
remainder property follow by scanning $0,\ldots,a$. The Euclidean
remainder decreases at every active step, so $a+b+1$ steps suffice for a
gcd. At each step retain the two linear combinations of the original
inputs. The identities $\gcd(u,v)=\gcd(v,u-qv)$ and
$u-qv$ equal to the updated linear combination prove the invariant.
This also proves Bezout's identity. For negative inputs apply the positive
algorithm to absolute values and adjust the signs of its coefficients.

For factorization, scan possible divisors $d=2,\ldots,a$. At a given
$d$, divide out at most $a$ copies. The invariant states that the product
of extracted factors times the current residual is $a$, and no previously
processed $d$ divides the residual. Whenever a divisor is first extracted
it is prime: a proper prime divisor would already have been removed.
That assertion is proved using the bounded search for the least divisor
of $d$ greater than one; the least divisor is prime because a nontrivial
factor of it would be smaller. Every assertion quantifies only over
integers bounded by $a$. Thus the finite invariant can be written with
characteristic functions and proved by quantifier-free induction.

If $\ell\nmid u$, primality gives $\gcd(\ell,u)=1$ by checking the
possible divisors of $\ell$; Bezout then shows that
$\ell\mid uv$ implies $\ell\mid v$. Repeated division by $\ell$, with
clock $|u|+|v|+1$, proves the valuation identity. Applied to the
factorizations already computed, it proves uniqueness of all prime
exponents.

In a coprime product, a prime occurs in at most one factor. Its valuation
in that factor therefore equals its valuation in the product, which is
divisible by $m$. Divide every exponent by $m$ and multiply the resulting
prime powers to obtain the root. For odd $m$ restore the original sign.
For $ab=2u^2$, all odd-prime valuations are even and the valuation at $2$
is odd in exactly the even factor. This gives the last assertion.
Each extracted root is bounded by the absolute value of its input; it may
alternatively be found by testing every candidate up to that bound.
\end{proof}

All later occurrences of ``coprime'' can be checked by this gcd algorithm.
All uses of valuation concern nonzero inputs and a fixed prime, or a
bounded list of prime factors of a displayed integer. Assertions such as
``every prime divisor has this property'' can therefore be read as bounded
assertions about that list.

\begin{lemma}[Primitive reduction]\label{lem:primitive}
For $n\geq1$, a positive solution of $x^n+y^n=z^n$ gives a positive
solution with pairwise coprime coordinates. Likewise, a positive solution
of $x^4+y^4=z^2$ gives one with $\gcd(x,y)=1$.
All divisions are bounded and the relevant right-hand coordinate does not
increase.
\end{lemma}

\begin{proof}
For the first equation take $d=\gcd(x,y)$. Prime valuations show
$d\mid z$: if $\ell^e\mid d$, then $\ell^{ne}\mid z^n$, so
$\ell^e\mid z$. Divide all three coordinates by $d$ and cancel $d^n$.
The new first two coordinates are coprime. A common prime of either of
them and the third divides the remaining one by the equation, a
contradiction. For the second equation the same argument gives
$d^2\mid z$; divide $x,y,z$ by $d,d,d^2$, respectively. The gcd algorithm
and the power-divisibility facts in Lemma~\ref{lem:integer-tools} give all
these witnesses constructively.
\end{proof}

